\lyseis{6}{Πρώτοι αριθμοί και ανάλυση: τα δομικά υλικά}
\sol{1} 12: 1, 2, 3, 4, 6, 12. \quad 18: 1, 2, 3, 6, 9, 18. \quad 30: 1, 2, 3, 5, 6, 10, 15, 30. \quad 36: 1, 2, 3, 4, 6, 9, 12, 18, 36. \quad 45: 1, 3, 5, 9, 15, 45.
\sol{2} Του 4: 0, 4, 8, 12, 16, 20. Του 7: 0, 7, 14, 21, 28, 35. Του 12: 0, 12, 24, 36, 48, 60.
\sol{3} Πρώτοι: 2, 13, 29, 41, 61, 89. Το 1 δεν είναι ούτε πρώτος ούτε σύνθετος. Σύνθετοι: 9 (3), 21 (3), 33 (3), 51 (3), 57 (3), 87 (3). Όλοι έχουν διαιρέτη το 3, γιατί το άθροισμα των ψηφίων τους (9, 3, 6, 6, 12, 15) διαιρείται με το 3.
\sol{4} 53, 59, 61, 67, 71, 73, 79.
\sol{5} α) $24=2^3\cdot 3$ \quad β) $60=2^2\cdot 3\cdot 5$ \quad γ) $72=2^3\cdot 3^2$ \quad δ) $100=2^2\cdot 5^2$ \quad ε) $150=2\cdot 3\cdot 5^2$ \quad στ) $504=2^3\cdot 3^2\cdot 7$.
\sol{6} α) $8\cdot 5=40$ \quad β) $2\cdot 9\cdot 7=126$ \quad γ) $4\cdot 3\cdot 25=300$ \quad δ) $27\cdot 11=297$.
\sol{7} α) Λάθος: ο 2 είναι άρτιος πρώτος. β) Λάθος. γ) Λάθος: $91=7\cdot 13$. δ) Σωστό. ε) Σωστό: $0=\alpha\cdot 0$. στ) Λάθος: $51=3\cdot 17$.
\sol{8} Ο 221 δεν διαιρείται με το 2, το 3 ($2+2+1=5$), το 5, το 7 ($7\cdot 31=217$) ούτε το 11 ($11\cdot 20=220$). Με το 13: $13\cdot 17=221$. Οι πρώτοι είναι 13 και 17.
\sol{9} Οι διαιρέτες του 48 είναι 1, 2, 3, 4, 6, 8, 12, 16, 24, 48: σακουλάκια με όποιον από αυτούς τους αριθμούς καραμελών. Το 47 είναι πρώτος, άρα έχει μόνο τους διαιρέτες 1 και 47: μόνο δύο τρόποι (47 σακουλάκια με 1 καραμέλα ή ένα σακουλάκι με 47).
\sol{10} Οι διαιρέτες του 60 είναι 1, 2, 3, 4, 5, 6, 10, 12, 15, 20, 30, 60. Αφαιρώ τον 1 και τον 60: μένουν 10, άρα 10 τρόποι (2, 3, 4, 5, 6, 10, 12, 15, 20 ή 30 παιδιά σε κάθε ομάδα). Το 61 είναι πρώτος, άρα κανένας τρόπος.
\sol{11} Οι πρώτοι ανάμεσα στο 20 και στο 30 είναι οι 23 και 29. Άθροισμα ψηφίων: $2+3=5$ και $2+9=11$. Ο αριθμός είναι ο 29.
\sol{12} Οι πρώτοι μέχρι το 60 είναι 17: 2, 3, 5, 7, 11, 13, 17, 19, 23, 29, 31, 37, 41, 43, 47, 53, 59. Το 2 σβήνει 29 αριθμούς (4, 6, 8, \dots, 60). Το 3 σβήνει ακόμα 9 αριθμούς: 9, 15, 21, 27, 33, 39, 45, 51, 57 (τα άλλα πολλαπλάσιά του είναι ήδη σβησμένα, γιατί είναι άρτια).
\sol{13} Για το 6: $1+2+3=6$. Για το 28: $1+2+4+7+14=28$. Σε κάθε περίπτωση το άθροισμα των διαιρετών (εκτός από τον ίδιο) ισούται με τον αριθμό.
\sol{14} $(3,\,5)$, $(5,\,7)$, $(11,\,13)$, $(17,\,19)$, $(29,\,31)$, $(41,\,43)$.
\sol{15} Οι αριθμοί είναι $1, 3, 9, 2, 6, 18, 4, 12, 36, 8, 24, 72$, δηλαδή 12 αριθμοί. Είναι όλοι οι διαιρέτες του 72.
\sol{16} Αν το 1 ήταν πρώτος, το 6 θα είχε πολλές αναλύσεις: $2\cdot 3$, $1\cdot 2\cdot 3$, $1\cdot 1\cdot 2\cdot 3$, \dots, όσες θέλεις. Η ανάλυση δεν θα ήταν μοναδική.
\sol{17} $101\cdot 103=10.100+303=10.403$. Αν δοκιμάζαμε τους πρώτους με τη σειρά (2, 3, 5, \dots, 97), θα αποτυγχάναμε 25 φορές, και μόνο ο 101 θα πετύχαινε. Με ακόμα μεγαλύτερους αριθμούς η δουλειά γίνεται τεράστια, και σε αυτό στηρίζεται η ασφάλεια των κωδικών.
